\section{A projective finite-type transfer with its exact positive-multiplicity gate}
\label{sec:cubic-projective-type-transfer}

This section is a supporting result with explicit arithmetic premises.  It does not assert
that an ordinary ramified generic automorphic representation with the
required residual datum has been supplied.  In particular, the generic
output of Flicker's split-prime Theorem~0.5 is not substituted for a
ramified output.

Let \(\mathcal O\) be the ring of integers of a finite extension of
\(\mathbb Q_p\), with fraction field \(E\) and residue field \(k\).
Enlarging \(E\) to split the finite groups in question is harmless.
For a finite group \(D\), a finite projective \(\mathcal O[D]\)-module
\(P\), and an irreducible characteristic-zero representation \(\sigma\),
write
\[
 m_\sigma(P)=
 \dim_{\overline E}\operatorname{Hom}_{\overline E[D]}
       (\sigma,P\otimes_{\mathcal O}\overline E).
 \tag{G72.1}
\]

\begin{lemma}[Projective characters only test the prime-to-\(p\) classes]
\label{G72-projective-character}
Suppose three characteristic-zero characters of \(D\) satisfy
\(\chi_\tau=\chi_s-\chi_t\) on every element of order prime to \(p\).
Then every finite projective \(\mathcal O[D]\)-module satisfies
\[
                  m_\tau(P)=m_s(P)-m_t(P).
 \tag{G72.2}
\]
The formula also holds for virtual characters, by linearity.
\end{lemma}
\begin{proof}
We first prove that \(\chi_P(g)=0\) whenever the order of \(g\)
is divisible by \(p\).  Write \(g=uh=hu\), where \(u\ne1\)
has \(p\)-power order and \(h\) has order prime to \(p\).
Restriction to \(\langle g\rangle\) preserves projectivity:
\(\mathcal O[D]\) is free over this subgroup ring, and a
projective module is a direct summand of a finite free module.
After a finite unramified coefficient extension containing the
\(\operatorname{ord}(h)\)-th roots of unity, the integral idempotents
of \(\langle h\rangle\) decompose the restriction into its
\(h\)-eigenspaces.  Each eigenspace is projective over
\(\mathcal O[\langle u\rangle]\).

This last commutative ring is local.  Indeed its reduction modulo the
uniformizer is \(k[X]/(X^{p^a}-1)=k[X]/((X-1)^{p^a})\), and every
maximal ideal of a finite algebra over the complete local ring
\(\mathcal O\) lies over its maximal ideal.  A finite projective
module over this local ring is free.  The trace of the nonidentity
element \(u\) on its regular representation is zero, so the trace
of \(uh\) on each \(h\)-eigenspace is zero.  Summing proves the
claim, and extension of coefficients does not change the original
trace.

Now apply the ordinary character inner product to \(\chi_P\)
and \(\chi_\tau-\chi_s+\chi_t\).  On the prime-to-\(p\)
classes the second factor is zero; on all other classes the first
factor is zero.  The inner product is therefore zero, giving
(G72.2).  No assertion about which simple modular constituents
occur in \(P/\mathfrak mP\) is needed.
\end{proof}

\begin{lemma}[The exact integral projectivity at a definite level]
\label{G72-definite}
Let \(U/\mathbb Q\) be a unitary group which is compact at infinity
and split at a rational prime \(q\ne p\).  Put
\(K_q=\mathrm{GL}_3(\mathbb Z_q)\),
\(K_q(1)=\ker(K_q\to\mathrm{GL}_3(\mathbb F_q))\), and
\(D_q=K_q/K_q(1)\).  Fix a finite free integral coefficient
module for the algebraic weight and the prescribed finite types
away from \(q\).  Suppose the compact level away from \(q\)
has a torsion-free factor at one auxiliary finite place.
Then the full lattice \(M\) of definite algebraic modular forms
at level \(K^qK_q(1)\) is free over \(\mathcal O[D_q]\).

Every direct summand selected by commuting integral ordinary
idempotents, integral type idempotents, and a local factor of a
finite integral away-Hecke algebra is projective over
\(\mathcal O[D_q]\).  The assertion requires the operators to
preserve the displayed integral lattice; it does not identify an
arbitrarily normalized operator with the source's ordinary operator.
\end{lemma}
\begin{proof}
\proofstep{The deck action and the coefficient fibers.}
The double-coset set
\[
 X=U(\mathbb Q)\backslash U(\mathbb A_f)/(K^qK_q(1))
 \tag{G72.3}
\]
is finite, as in the definite-form construction in Flicker,
\S5.  Its right \(D_q\)-action is free.  To see this, a
stabilizer of a point represented by \(x\) gives a rational
element in \(U(\mathbb Q)\cap x(K^qK_q)x^{-1}\).
This arithmetic stabilizer is finite: it is a discrete subgroup
of the compact set obtained by adding the compact infinite factor.
Its image at the torsion-free auxiliary factor is trivial, so
the rational element is the identity.  The alleged deck element
then lies in \(K_q(1)\), as required.

Evaluation realizes forms as sections of a \(D_q\)-equivariant
bundle of finite free \(\mathcal O\)-modules over \(X\).
The coefficient action, including the algebraic action at \(p\),
is built into these fibers.  Choose one point in each free orbit
and a basis of its fiber; translate that basis around the orbit.
There is no stabilizer relation to impose.  The contribution of
each orbit is \(\mathcal O[D_q]\) tensored with its chosen free
fiber.  This proves freeness, without division by \(|D_q|\).
This is the equivariant version of the evaluation formula preceding
Flicker's Lemma~5.6; the torsion-free-factor observation is made
explicitly in his \S6 immediately after Definition~6.1.

\proofstep{The integral projectors.}
If an integral endomorphism \(T\) acts on a finite free
\(\mathcal O\)-module, then \(T^{n!}\) has a limit in its
endomorphism ring.  For every \(a\), the module modulo
\(\mathfrak m^a\) is finite.  Its Fitting decomposition for
\(T\) is a direct sum of a nilpotent part and a part on which
\(T\) is invertible.  For sufficiently large \(n\),
\(T^{n!}\) is zero on the first and the identity on the second.
These idempotents are compatible as \(a\) varies.  Completeness
therefore gives an integral idempotent \(e_T\).  Over a splitting
field its image is exactly the sum of the generalized eigenspaces
whose eigenvalues are units.  This also follows by the coprime
unit-root/nonunit-root factorization of the characteristic
polynomial.  Commuting operators give commuting limits.

A finite commutative \(\mathcal O\)-algebra of integral Hecke
operators is complete.  The idempotents of its Artinian reduction
lift uniquely through the successive powers of \(\mathfrak m\),
since the derivative of \(X^2-X\) is a unit at either idempotent.
Thus its local factors are cut out by integral idempotents as well.
All these idempotents commute with \(D_q\) when the original
operators do.  Their product selects a direct summand of the free
\(\mathcal O[D_q]\)-module \(M\), proving projectivity.
This proof does not use semisimplicity modulo \(p\).
\end{proof}

\begin{remark}[Removing an auxiliary neat level at prime-to-\(p\) cost]
For the actual fixed field \(K=\mathbb Q(\omega)\) and \(p>3\),
one may impose the torsion-free factor without retaining a new
away type.  Choose, outside any fixed finite set, a prime \(r\)
with \(r\equiv1\pmod3\) and whose image modulo \(p\) has
order \(p-1\).  The Chinese remainder theorem and Dirichlet's
theorem give such primes.  The group is split at \(r\), and
\[
 |\mathrm{GL}_3(\mathbb F_r)|
       =r^3(r-1)(r^2-1)(r^3-1)
       \quad\hbox{is prime to }p.
 \tag{G72.4}
\]
For \(r\ge3\), the principal congruence group
\(1+rM_3(\mathbb Z_r)\) is torsion free: its convergent
matrix logarithm is injective, and the logarithm of a finite power
of one element is the corresponding integer multiple of its
logarithm.  After working at this principal level, averaging over
\(\mathrm{GL}_3(\mathbb F_r)\) is an integral idempotent.
Its image is the original hyperspecial level at \(r\).
It commutes with the \(q\)-deck action and with ordinary and
Hecke operators at other places.  Subsequent integral localization
at the restored spherical Hecke algebra also preserves
projectivity.  Thus no new ramification at \(r\) is asserted or
needed.  For a different CM base, the existence of an auxiliary
place with these simultaneous properties is a separate hypothesis.
\end{remark}

\begin{theorem}[AC-G72: the exact positive-multiplicity finite-type bridge]
\label{AC-G72}
Let \(p>3\), let \(q\) be an odd prime with
\(q\equiv-1\pmod p\), and
\(D_q=\mathrm{GL}_3(\mathbb F_q)\).
Let \(\tau=\tau_\theta\) be the irreducible finite representation
constructed in AC-G72-character from a nontrivial \(p\)-power
character of the nonsplit rank-two torus.  That proposition proves
the precise finite-group character identity
\[
       \chi_\tau(g)=\chi_{\mathrm{St}_3}(g)-1
       \qquad(g\in D_q\text{ of order prime to }p).
 \tag{G72.5}
\]
Let \(P\) be the actual ordinary, away-type and residual-Hecke
summand of the definite-form lattice from Lemma~\ref{G72-definite},
with all its stated integral hypotheses.  Then
\[
             m_\tau(P)=m_{\mathrm{St}_3}(P)-m_1(P).
 \tag{G72.6}
\]
In particular, an actual \(\tau\)-eigensystem in this same
weight, ordinary summand and residual away-Hecke block exists if
\[
                     m_{\mathrm{St}_3}(P)>m_1(P).
 \tag{G72.7}
\]

Here is a sufficient automorphic condition for (G72.7).  In the
characteristic-zero automorphic decomposition of \(P\), suppose
every spherical \(q\)-component is an irreducible unramified
principal series, and at least one component of positive
multiplicity is generic, Iwahori spherical and nonspherical.
Then (G72.7) holds.  The nonspherical generic component may be
of type IIa or IIIa in Flicker's \S11.5, Table~A; it need not
be a full smooth Steinberg representation.
\end{theorem}
\begin{proof}
Projectivity and (G72.5) give (G72.6) by
Lemma~\ref{G72-projective-character}.  If its right side is positive,
\(\operatorname{Hom}_{D_q}(\tau,P\otimes\overline E)\)
is a nonzero finite-dimensional space, stable under the commuting
away-Hecke algebra and ordinary projectors.  A commuting algebra
over an algebraically closed field has a common eigenvector on
a nonzero finite-dimensional invariant space: choose a minimal
nonzero invariant subspace, on which every operator is scalar.
This gives an actual Hecke eigensystem in the selected space.
Its away eigenvalues reduce to the specified maximal ideal
because the space was selected by that integral local factor.
No modular Jordan--H\"older constituent is lifted by fiat.

For completeness, we compute the multiplicity difference from
the local representations.  Let \(R\) be the nontrivial
constituent of the permutation representation on the projective
plane over \(\mathbb F_q\).  The action on points is doubly
transitive, so its commuting algebra has dimension two; hence
the permutation module is \(1\oplus R\) with \(R\)
irreducible.  The complete-flag permutation module has commuting
algebra of dimension six, by the six relative positions of
two flags.  The space of maps from the point-permutation module
to it has dimension three, by the three positions of a point
relative to a fixed flag.  Thus the complete-flag module contains
\(1\) once and \(R\) twice.  The remaining character has
norm \(6-1-4=1\) and degree \(q^3\); equivalently it is
the top reduced homology character of the rank-two building.
It is the finite Steinberg character.  Therefore
\[
 \operatorname{Ind}_{B(\mathbb F_q)}^{D_q}1
       =1\oplus R^{\oplus2}\oplus\mathrm{St}_3,
 \qquad
 R_{\mathrm{GL}_2\times\mathrm{GL}_1}^{D_q}
       (\mathrm{St}_2\otimes1)=R\oplus\mathrm{St}_3.
 \tag{G72.8}
\]
The second equality follows from transitivity of induction and
\(\operatorname{Ind}_{B(\mathbb F_q)}^{\mathrm{GL}_2(\mathbb F_q)}1
=1\oplus\mathrm{St}_2\): induction of the trivial Levi
representation is the point-permutation module.

For an irreducible unramified principal series of
\(\mathrm{GL}_3(\mathbb Q_q)\), Iwasawa decomposition identifies
its \(K_q(1)\)-invariants with the first module in (G72.8).
Consequently its finite Steinberg and trivial multiplicities
are both one.  A type IIa representation is an irreducible
parabolic induction of an unramified twist of the smooth
\(\mathrm{GL}_2\) Steinberg representation and an unramified
character.  Its \(K_q(1)\)-invariants are the second module
in (G72.8), so these two multiplicities are one and zero.
Here taking invariants is exact in characteristic zero, and
the smooth rank-one Steinberg quotient restricts to the finite
rank-one Steinberg quotient.  For type IIIa the same standard
parabolic quotient defining smooth Steinberg gives the finite
\(\mathrm{St}_3\) on \(K_q(1)\)-invariants.  Its
multiplicities are again one and zero.

Every other nonspherical component contributes zero to the
trivial finite-type multiplicity and a nonnegative integer to
the Steinberg multiplicity.  Indeed finite trivial invariants
are precisely \(K_q\)-invariants.  Summing these differences
with their nonnegative automorphic and ordinary multiplicities
proves strict positivity.  Flicker's complete Table~A says
that the generic Iwahori-spherical possibilities are exactly
I, IIa and IIIa; type I is spherical.  Thus the stated
nonspherical generic hypothesis is precisely sufficient.
\end{proof}

\begin{remark}[Why Steinberg occurrence and projectivity alone are insufficient]
Take \(P=\mathcal O[D_q/B(\mathbb F_q)]\).  Since
\(|B(\mathbb F_q)|=q^3(q-1)^3\) is prime to \(p\), its
trivial integral representation is projective, by averaging,
and induction makes \(P\) projective over \(\mathcal O[D_q]\).
Equation (G72.8) gives \(m_{\mathrm{St}_3}(P)=m_1(P)=1\).
Thus (G72.6) gives \(m_\tau(P)=0\), although a Steinberg
finite type is present.  The removed trivial finite type means
sphericity at \(q\); it does not mean that the global
automorphic representation is one-dimensional.  A global
nonabelian hypothesis does not eliminate this counterexample
to the fragmentary finite-module inference.
\end{remark}

\paragraph{Exact primary-source scope and remaining arithmetic gates.}
Lang--Pollack--Wake, arXiv:2609.05144v1, \S3.2, Lemmas~3.6--3.8
and the proof of Theorem~3.1(2a), provide a precise rank-two
projective-module precedent.  Their theorem is about the stated
rank-two finite group and is not cited as a rank-three
automorphic theorem.  The character-vanishing and definite-form
projectivity arguments above are proved here.  The intended
rank-three type comes from a regular \(p\)-power character of
\(\mathbb F_{q^2}^{\times}\) and the \((2,1)\) Levi.
Its irreducibility and (G72.5) are proved by the complete elementary
calculation in AC-G72-character.  Its exact smooth local-type
interpretation is a separate depth-zero-type interface; the
present theorem only produces an actual finite-type eigensystem.

Flicker, \emph{On level-raising congruences}, Theorem~0.5 and
its complete proof in \S11.5, supply a type I or IIa output,
not a forced type IIa output.  A type I contribution has
zero difference in (G72.7).  Thus that theorem alone has not
been shown to supply the strict inequality used here.
Matching its proof to the actual integral ordinary projector
and proving the required residual-block spherical condition
also remain explicit arithmetic tasks.  The existence of a
finite type vector alone does not assert attached Galois
local-global compatibility, the exact type at every other
away prime, the correct ordered local component at \(p\),
or the actual curve-support containment \(J\subseteq\ker f\).

The proposed auxiliary conditions \(q\equiv-1\pmod p\) and
\(\gamma(q)=-1\) are not asserted available for every literal
source datum.  The source parity condition can obstruct them.
Replacing a rational character by its product with an even
quadratic character which becomes trivial on a chosen host
is a possible host-level comparison, not a scalar twist of
an arbitrary irreducible rank-three representation and not an
already proved source-frame transport.

\paragraph{Geometric dictionary.}
The free finite deck action and coefficient bundle express
descent of sections along a finite cover, in the already
audited SGA~1, Expos\'e~V, \S\S4 and 7 framework.  The
projective integral lattice records a direct summand of that
descent module.  This dictionary supplies neither a missing
automorphic seed nor the strict multiplicity inequality.
